\section{Потоки в сетях. Теорема Форда и Фалкерсона. Алгоритмы нахождения максимального потока}
\label{sec:network-flows}

\subsection{Транспортная сеть и допустимый поток}

Транспортной сетью называется ориентированный граф
\[
    G=(V,E)
\]
с выделенными вершинами $s$ (исток) и $t$ (сток) и пропускными способностями
\[
    c\colon E\to\mathbb R_{\ge0}.
\]

Поток --- функция $f\colon E\to\mathbb R_{\ge0}$, удовлетворяющая:
\begin{enumerate}
 \item ограничению пропускной способности
 \[
    0\le f(u,v)\le c(u,v);
 \]
 \item закону сохранения во внутренних вершинах
 \[
    \sum_{(u,v)\in E}f(u,v)
    =
    \sum_{(v,w)\in E}f(v,w),
    \qquad v\in V\setminus\{s,t\}.
 \]
\end{enumerate}

Величина потока --- чистый выход из источника:
\[
 |f|=
 \sum_{(s,v)\in E}f(s,v)-
 \sum_{(v,s)\in E}f(v,s).
\]
Из закона сохранения следует, что эта величина равна чистому входу в сток.

Удобна также антисимметричная запись потока на всех упорядоченных парах:
$f(u,v)=-f(v,u)$. Тогда закон сохранения имеет компактный вид
\[
    \sum_{v\in V} f(u,v)=0,
    \qquad u\ne s,t.
\]

\subsection{Остаточная сеть}

Идея алгоритмов максимального потока состоит в том, что уже проведённый поток
можно не только увеличивать, но и частично отменять. Для каждой исходной дуги
$(u,v)$ создаются остаточные возможности:
\[
    c_f^{+}(u,v)=c(u,v)-f(u,v),
    \qquad
    c_f^{-}(v,u)=f(u,v).
\]
Остаточная сеть $G_f$ содержит дуги положительной остаточной пропускной
способности. Если в исходной сети одновременно имеются $(u,v)$ и $(v,u)$, их
остаточные вклады нужно учитывать совместно (или хранить как разные объекты
рёбер); это устраняет неоднозначность простой записи ``прямая/обратная дуга''.

Путь $P$ из $s$ в $t$ в остаточной сети называется \emph{увеличивающим}. Его
пропускная способность
\[
    \Delta(P)=\min_{e\in P} c_f(e).
\]
Увеличение потока на $\Delta$ по прямой остаточной дуге увеличивает исходный
поток, а по обратной --- уменьшает его. После такой операции поток остаётся
допустимым и его величина возрастает на $\Delta$.

\subsection{Разрез и слабая двойственность}

$s$--$t$-разрез --- разбиение
\[
 V=S\sqcup T,
 \qquad s\in S,\quad t\in T.
\]
Его пропускная способность
\[
    c(S,T)=\sum_{u\in S,\,v\in T}c(u,v).
\]

Для любого потока справедливо тождество потока через разрез:
\[
 |f|=
 \sum_{u\in S,v\in T}f(u,v)
 -
 \sum_{u\in T,v\in S}f(u,v).
\]
Внутренние потоки внутри $S$ взаимно сокращаются благодаря закону сохранения.
Отсюда немедленно
\[
    |f|\le c(S,T).
\]
Это \textbf{слабая двойственность}: любой разрез даёт верхнюю оценку любого
допустимого потока.

\subsection{Критерий максимальности через увеличивающие пути}

\textbf{Теорема.} Допустимый поток $f$ максимален тогда и только тогда, когда в
остаточной сети $G_f$ нет увеличивающего пути из $s$ в $t$.

Необходимость очевидна: если путь есть, поток можно увеличить. Для достаточности
пусть увеличивающего пути нет и $S$ --- множество вершин, достижимых из $s$ в
остаточной сети. Тогда $t\notin S$. Для любой исходной дуги из $S$ в $T$ её
остаточная ёмкость нулевая, то есть дуга насыщена. Для любой исходной дуги из
$T$ в $S$ поток равен нулю, иначе существовала бы положительная обратная
остаточная дуга из $S$ в $T$. Поэтому
\[
    |f|=c(S,T).
\]
Но поток не может превосходить никакой разрез, следовательно, найденный поток
максимален.

\subsection{Теорема максимального потока и минимального разреза}

Из предыдущего критерия следует центральный результат.

\textbf{Теорема Форда--Фалкерсона (max-flow min-cut).}
\[
    \boxed{
    \max_f |f|=
    \min_{(S,T)} c(S,T)
    }.
\]
То есть оптимальное значение прямой задачи ``максимизировать поток'' равно
оптимальному значению двойственной комбинаторной задачи ``минимизировать
пропускную способность разреза''.

Теорема даёт не только значение оптимума, но и \emph{сертификат} оптимальности:
если предъявлены поток $f$ и разрез $(S,T)$ с
\[
    |f|=c(S,T),
\]
то оба автоматически оптимальны.

\subsection{Алгоритм Форда--Фалкерсона}

Алгоритм начинается с нулевого потока и повторяет:
\begin{enumerate}
 \item построить остаточную сеть;
 \item найти увеличивающий путь $P$;
 \item вычислить
 \[
    \Delta=\min_{e\in P}c_f(e);
 \]
 \item увеличить поток вдоль $P$ на $\Delta$.
\end{enumerate}
Если увеличивающего пути нет, по теореме поток максимален.

Если все пропускные способности целые, то начиная с целого потока алгоритм
сохраняет целочисленность. Каждое увеличение повышает $|f|$ как минимум на $1$,
поэтому при максимальном потоке величины $F$
\[
    T=O(mF)
\]
для реализации, где один поиск пути занимает $O(m)$. Эта оценка
псевдополиномиальна: зависит от численного значения $F$, а не только от длины
входа.

Для произвольных вещественных ёмкостей выбор плохих путей может приводить к
очень долгой работе, а при иррациональных значениях классическая схема даже
может не завершиться. Поэтому практически используют специальные правила выбора
увеличивающих путей.

\subsection{Алгоритм Эдмондса--Карпа}

Эдмондс--Карп выбирает кратчайший по числу дуг увеличивающий путь с помощью BFS.
Его сложность
\[
    O(nm^2).
\]
Идея оценки: расстояния от $s$ в остаточной сети не уменьшаются, а каждая дуга
может становиться критической на кратчайшем увеличивающем пути лишь $O(n)$ раз.
Каждый BFS стоит $O(m)$.

Главное преимущество --- полиномиальная оценка, не зависящая от значений
пропускных способностей.

\subsection{Алгоритм Диница}

Диниц на каждой фазе:
\begin{enumerate}
 \item BFS строит уровни $level(v)$ и слоистую остаточную сеть;
 \item в этой сети DFS находит \emph{блокирующий поток}, после которого ни один
       путь $s$--$t$, использующий только дуги между соседними уровнями, больше
       нельзя увеличить.
\end{enumerate}
После каждой фазы расстояние от $s$ до $t$ увеличивается, поэтому фаз не более
$n-1$. Для общего графа стандартная оценка
\[
    O(n^2m).
\]
На специальных классах сетей возможны лучшие оценки.

\subsection{Теорема о целочисленности}

Если все пропускные способности целые, существует максимальный поток, целый на
каждой дуге. Это следует из Форда--Фалкерсона: начиная с нулевого целого потока,
каждое $\Delta$ является целым.

Это свойство принципиально для дискретных приложений: единица потока может
интерпретироваться как выбор ребра, назначение задачи, перевозка одного объекта и
т.~п.

\subsection{Сведение паросочетания к потоку}

Пусть дан двудольный граф $G=(L\sqcup R,E)$. Строим сеть:
\begin{itemize}
 \item $s\to L$ с ёмкостями $1$;
 \item каждое ребро $L\to R$ с ёмкостью $1$;
 \item $R\to t$ с ёмкостями $1$.
\end{itemize}
Целочисленный поток величины $k$ соответствует паросочетанию из $k$ рёбер, и
обратно. Поэтому максимальное паросочетание находится как максимальный поток.
Аналогично через поток формулируются задачи о непересекающихся путях и некоторые
задачи распределения ресурсов.

\subsection{Что важно сказать на экзамене}

Нужно последовательно определить допустимый поток, остаточную сеть, увеличивающий
путь и разрез. Далее доказать оценку $|f|\le c(S,T)$ и сформулировать критерий:
поток максимален тогда и только тогда, когда увеличивающих путей нет. Отсюда
получается теорема max-flow min-cut. Затем сравнить Форд--Фалкерсона,
Эдмондса--Карпа и Диница и упомянуть целочисленность максимального потока.

\newpage
